% ============================================================
%  CHAPITRE : LIMITES DE FONCTIONS
%  Mazava Maths — Terminale S
%  Sans logarithme ni exponentielle dans les exemples
% ============================================================

\chapter{Limites de fonctions}

\begin{objectifs}
	\begin{itemize}
		\item Calculer la limite d'une fonction en un point ou à l'infini.
		\item Identifier et lever les formes indéterminées.
		\item Utiliser les théorèmes de comparaison et des gendarmes.
		\item Étudier les branches infinies d'une courbe.
	\end{itemize}
\end{objectifs}

% ----------------------------------------------------------------
\section{Notion de limite}
% ----------------------------------------------------------------

\begin{definition}
	Soit $f$ une fonction et $\ell \in \R$.
	
	\begin{itemize}
		\item \textbf{Limite finie en l'infini :}
		$\displaystyle\lim_{x\to+\infty}f(x)=\ell$ signifie que $f(x)$
		se rapproche de $\ell$ quand $x$ devient très grand.
		
		\item \textbf{Limite infinie en l'infini :}
		$\displaystyle\lim_{x\to+\infty}f(x)=+\infty$ signifie que $f(x)$
		devient aussi grand qu'on veut.
		
		\item \textbf{Limite en un point :}
		$\displaystyle\lim_{x\to a}f(x)=\ell$ signifie que $f(x)$
		se rapproche de $\ell$ quand $x$ se rapproche de $a$
		(sans nécessairement atteindre $a$).
		
		\item \textbf{Limite à gauche / à droite :}
		On note $\displaystyle\lim_{x\to a^-}f(x)$ (par valeurs inférieures)
		et $\displaystyle\lim_{x\to a^+}f(x)$ (par valeurs supérieures).
	\end{itemize}
\end{definition}

\begin{remarque}
	Si $\displaystyle\lim_{x\to a^-}f(x) \neq \lim_{x\to a^+}f(x)$,
	alors $f$ n'a \textbf{pas de limite} en $a$.
\end{remarque}

% ----------------------------------------------------------------
\section{Limites des fonctions de référence}
% ----------------------------------------------------------------

\begin{propriete}{Limites usuelles à connaître}{}
	\begin{center}
		\renewcommand{\arraystretch}{2.0}
		\begin{tabular}{|c|c|c|}
			\hline
			\textbf{Fonction} & \textbf{Limite} & \textbf{Résultat} \\
			\hline
			$x^n$ ($n\geq 1$) & $x\to+\infty$ & $+\infty$ \\
			\hline
			$x^n$ ($n$ pair) & $x\to-\infty$ & $+\infty$ \\
			\hline
			$x^n$ ($n$ impair) & $x\to-\infty$ & $-\infty$ \\
			\hline
			$\dfrac{1}{x}$ & $x\to 0^+$ & $+\infty$ \\
			\hline
			$\dfrac{1}{x}$ & $x\to 0^-$ & $-\infty$ \\
			\hline
			$\dfrac{1}{x^n}$ ($n\geq 1$) & $x\to\pm\infty$ & $0$ \\
			\hline
			$\sqrt{x}$ & $x\to+\infty$ & $+\infty$ \\
			\hline
			$\sqrt{x}$ & $x\to 0^+$ & $0$ \\
			\hline
			$\cos x$, $\sin x$ & $x\to\pm\infty$ & pas de limite \\
			\hline
		\end{tabular}
	\end{center}
\end{propriete}

% ----------------------------------------------------------------
\section{Opérations sur les limites}
% ----------------------------------------------------------------

\begin{propriete}{Règles opératoires}{}
	Si $\displaystyle\lim_{x\to a}f(x)=\ell$ et $\displaystyle\lim_{x\to a}g(x)=m$
	($\ell, m \in \R$ ou $\pm\infty$), alors :
	
	\begin{center}
		\renewcommand{\arraystretch}{2.0}
		\begin{tabular}{|c|c|c|}
			\hline
			\textbf{Opération} & \textbf{Résultat} & \textbf{Cas indéterminé} \\
			\hline
			$f + g$ & $\ell + m$ & $+\infty - \infty$ \quad \textbf{FI} \\
			\hline
			$f \times g$ & $\ell \times m$ & $0 \times \infty$ \quad \textbf{FI} \\
			\hline
			$\dfrac{f}{g}$ & $\dfrac{\ell}{m}$ (si $m\neq 0$)
			& $\dfrac{0}{0}$ et $\dfrac{\infty}{\infty}$ \quad \textbf{FI} \\
			\hline
			$f^n$ & $\ell^n$ & — \\
			\hline
			$\sqrt{f}$ & $\sqrt{\ell}$ (si $\ell\geq 0$) & — \\
			\hline
		\end{tabular}
	\end{center}
	
	\medskip
	\textbf{Règles pour les infinis (non indéterminés) :}
	\begin{center}
		\renewcommand{\arraystretch}{1.8}
		\begin{tabular}{|c|c|}
			\hline
			$\ell + \infty = +\infty$ & $\ell \times (+\infty) = +\infty$ si $\ell>0$ \\
			\hline
			$\ell - \infty = -\infty$ & $\ell \times (+\infty) = -\infty$ si $\ell<0$ \\
			\hline
			$\dfrac{\ell}{\pm\infty} = 0$ & $\dfrac{\ell}{0^+} = +\infty$ si $\ell>0$ \\
			\hline
			$+\infty + \infty = +\infty$ & $\dfrac{\ell}{0^-} = -\infty$ si $\ell>0$ \\
			\hline
			$+\infty \times +\infty = +\infty$ & $\infty \times \infty = \infty$ \\
			\hline
		\end{tabular}
	\end{center}
\end{propriete}

% ----------------------------------------------------------------
\section{Formes indéterminées et méthodes de levée}
% ----------------------------------------------------------------

\begin{attention}
	Les \textbf{formes indéterminées} (FI) sont :
	\[
	\frac{0}{0}, \quad \frac{\infty}{\infty}, \quad
	+\infty - \infty, \quad 0 \times \infty.
	\]
	Elles nécessitent un \textbf{traitement spécifique} avant de conclure.
\end{attention}

% --------------------------------
\subsection{FI de type \texorpdfstring{$\dfrac{\infty}{\infty}$}{inf/inf} -- Fonctions rationnelles}
% --------------------------------

\begin{methode}
	Pour $\displaystyle\lim_{x\to\pm\infty}\frac{P(x)}{Q(x)}$
	où $P$ et $Q$ sont des polynômes :
	\begin{enumerate}
		\item \textbf{Factoriser} numérateur et dénominateur par
		le \textbf{terme de plus haut degré} de chacun,
		\textbf{ou} identifier directement les termes dominants.
		\item Simplifier.
		\item Conclure avec les limites usuelles.
	\end{enumerate}
	\textbf{Règle du coup d'œil :}
	$\displaystyle\lim_{x\to\pm\infty}\frac{a_nx^n+\cdots}{b_mx^m+\cdots}
	= \lim_{x\to\pm\infty}\frac{a_nx^n}{b_mx^m}$.
\end{methode}

\begin{exemple}
	Calculer $\displaystyle L_1=\lim_{x\to+\infty}\frac{3x^2-2x+1}{x^2+5x-3}$.
	
	On factorise par $x^2$ :
	\[
	L_1 = \lim_{x\to+\infty}
	\frac{x^2\!\left(3-\frac{2}{x}+\frac{1}{x^2}\right)}
	{x^2\!\left(1+\frac{5}{x}-\frac{3}{x^2}\right)}
	= \frac{3-0+0}{1+0-0} = 3.
	\]
\end{exemple}

\begin{exemple}
	Calculer $\displaystyle L_2=\lim_{x\to-\infty}\frac{2x^3-x}{x^2+1}$.
	
	Terme dominant : $\dfrac{2x^3}{x^2}=2x$ :
	\[
	L_2 = \lim_{x\to-\infty}2x = -\infty.
	\]
\end{exemple}

% --------------------------------
\subsection{FI de type \texorpdfstring{$\dfrac{0}{0}$}{0/0}}
% --------------------------------

\begin{methode}
	Pour une FI $\dfrac{0}{0}$ :
	\begin{enumerate}
		\item \textbf{Factoriser} numérateur et dénominateur
		(par $(x-a)$, par un facteur commun).
		\item Utiliser les \textbf{identités remarquables} :
		$a^2-b^2=(a-b)(a+b)$, $a^3-b^3=(a-b)(a^2+ab+b^2)$.
		\item \textbf{Simplifier} le facteur qui s'annule.
		\item Calculer la limite de l'expression simplifiée.
	\end{enumerate}
\end{methode}

\begin{exemple}
	Calculer $\displaystyle L_3=\lim_{x\to 2}\frac{x^2-4}{x-2}$.
	
	$\dfrac{x^2-4}{x-2} = \dfrac{(x-2)(x+2)}{x-2} = x+2$
	\quad (pour $x\neq 2$).
	
	Donc $L_3 = 2+2 = 4$.
\end{exemple}

\begin{exemple}
	Calculer $\displaystyle L_4=\lim_{x\to 1}\frac{x^3-1}{x^2-1}$.
	
	$x^3-1=(x-1)(x^2+x+1)$ et $x^2-1=(x-1)(x+1)$, donc :
	\[
	\frac{x^3-1}{x^2-1}
	= \frac{(x-1)(x^2+x+1)}{(x-1)(x+1)}
	= \frac{x^2+x+1}{x+1} \xrightarrow[x\to 1]{} \frac{3}{2}.
	\]
\end{exemple}

% --------------------------------
\subsection{FI de type \texorpdfstring{$+\infty - \infty$}{inf - inf}}
% --------------------------------

\begin{methode}
	Pour une FI $+\infty-\infty$ :
	\begin{itemize}
		\item \textbf{Si expression rationnelle :} réduire au même dénominateur.
		\item \textbf{Si racines carrées :} multiplier par la \textbf{quantité conjuguée}
		$\sqrt{A}+\sqrt{B}$ ou $f(x)+g(x)$.
		\item \textbf{Factoriser} par le terme dominant.
	\end{itemize}
\end{methode}

\begin{exemple}
	Calculer $\displaystyle L_5=\lim_{x\to+\infty}\left(\sqrt{x^2+3x}-x\right)$.
	
	On multiplie par la quantité conjuguée :
	\[
	\sqrt{x^2+3x}-x
	= \frac{(\sqrt{x^2+3x}-x)(\sqrt{x^2+3x}+x)}{\sqrt{x^2+3x}+x}
	= \frac{x^2+3x-x^2}{\sqrt{x^2+3x}+x}
	= \frac{3x}{\sqrt{x^2+3x}+x}.
	\]
	Pour $x>0$ : $\sqrt{x^2+3x}=x\sqrt{1+3/x}$, donc :
	\[
	\frac{3x}{x\sqrt{1+\frac{3}{x}}+x}
	= \frac{3}{\sqrt{1+\frac{3}{x}}+1}
	\xrightarrow[x\to+\infty]{} \frac{3}{1+1} = \frac{3}{2}.
	\]
\end{exemple}

\begin{exemple}
	Calculer $\displaystyle L_6=\lim_{x\to+\infty}(x^2-3x+1)$.
	
	On factorise par $x^2$ :
	$x^2-3x+1 = x^2\!\left(1-\dfrac{3}{x}+\dfrac{1}{x^2}\right)\to+\infty$.
\end{exemple}

% --------------------------------
\subsection{FI de type \texorpdfstring{$0\times\infty$}{0 x inf}}
% --------------------------------

\begin{methode}
	Pour une FI $0\times\infty$, transformer en $\dfrac{0}{0}$ ou
	$\dfrac{\infty}{\infty}$ :
	\[
	f(x)\cdot g(x) = \frac{f(x)}{1/g(x)} \quad \text{ou} \quad
	\frac{g(x)}{1/f(x)}.
	\]
\end{methode}

\begin{exemple}
	Calculer $\displaystyle L_7=\lim_{x\to 0^+} x\cdot\frac{1}{\sqrt{x}}
	= \lim_{x\to 0^+}\sqrt{x} = 0$.
	
	(On écrit $x\cdot\dfrac{1}{\sqrt{x}}=\dfrac{x}{\sqrt{x}}=\sqrt{x}$.)
\end{exemple}

% --------------------------------
\subsection{Limites et racines carrées en \texorpdfstring{$\dfrac{0}{0}$}{0/0}}
% --------------------------------

\begin{methode}
	Pour $\displaystyle\lim_{x\to a}\frac{\sqrt{f(x)}-\sqrt{g(x)}}{h(x)}$,
	multiplier par $\dfrac{\sqrt{f(x)}+\sqrt{g(x)}}{\sqrt{f(x)}+\sqrt{g(x)}}$
	pour éliminer les racines au numérateur.
\end{methode}

\begin{exemple}
	Calculer $\displaystyle L_8=\lim_{x\to 1}\frac{\sqrt{x}-1}{x-1}$.
	
	\[
	\frac{\sqrt{x}-1}{x-1}
	= \frac{\sqrt{x}-1}{(\sqrt{x}-1)(\sqrt{x}+1)}
	= \frac{1}{\sqrt{x}+1}
	\xrightarrow[x\to 1]{} \frac{1}{2}.
	\]
\end{exemple}

% ----------------------------------------------------------------
\section{Limite d'une fonction composée}
% ----------------------------------------------------------------

\begin{theoreme}{Limite de la composée}
	Si $\displaystyle\lim_{x\to a}g(x)=b$ et $\displaystyle\lim_{t\to b}f(t)=\ell$,
	alors :
	\[
	\lim_{x\to a} f\bigl(g(x)\bigr) = \ell.
	\]
\end{theoreme}

\begin{methode}
	\begin{enumerate}
		\item Identifier la fonction intérieure $u=g(x)$ et calculer $\lim u$.
		\item Calculer $\lim f(u)$ quand $u$ tend vers la valeur trouvée.
		\item Conclure.
	\end{enumerate}
\end{methode}

\begin{exemple}
	Calculer $\displaystyle\lim_{x\to+\infty}f(x)$ avec
	$f(x)=\dfrac{1}{x^2+1}$.
	
	$u=x^2+1\to+\infty$, puis $\dfrac{1}{u}\to 0$. Donc $L=0$.
\end{exemple}

% ----------------------------------------------------------------
\section{Théorèmes de comparaison}
% ----------------------------------------------------------------

\begin{theoreme}{Théorème de comparaison}
	Soit $f$ et $g$ deux fonctions telles que $f(x)\leq g(x)$
	au voisinage de $a$.
	\begin{itemize}
		\item Si $\displaystyle\lim_{x\to a}f(x)=+\infty$, alors
		$\displaystyle\lim_{x\to a}g(x)=+\infty$.
		\item Si $\displaystyle\lim_{x\to a}g(x)=-\infty$, alors
		$\displaystyle\lim_{x\to a}f(x)=-\infty$.
	\end{itemize}
\end{theoreme}

\begin{theoreme}{Théorème des gendarmes}
	Soit $f$, $g$, $h$ trois fonctions telles que,
	au voisinage de $a$ :
	\[
	g(x) \leq f(x) \leq h(x).
	\]
	Si $\displaystyle\lim_{x\to a}g(x)=\lim_{x\to a}h(x)=\ell$,
	alors $\displaystyle\lim_{x\to a}f(x)=\ell$.
\end{theoreme}

\begin{astucebac}
	Le théorème des gendarmes est très souvent utilisé au BAC
	avec des fonctions oscillantes du type $\dfrac{\sin x}{x}$
	ou $\dfrac{\cos x}{x^2}$ qui sont encadrées par $\pm\dfrac{1}{|x|}$.
\end{astucebac}

\begin{exemple}
	Montrer que $\displaystyle\lim_{x\to+\infty}\frac{\sin x}{x}=0$.
	
	Pour tout $x>0$ : $-1\leq\sin x\leq 1$, donc
	$\dfrac{-1}{x}\leq\dfrac{\sin x}{x}\leq\dfrac{1}{x}$.
	
	Or $\displaystyle\lim_{x\to+\infty}\frac{-1}{x}=0$ et
	$\displaystyle\lim_{x\to+\infty}\frac{1}{x}=0$.
	
	Par le théorème des gendarmes :
	$\displaystyle\lim_{x\to+\infty}\frac{\sin x}{x}=0$.
\end{exemple}

\begin{exemple}
	Soit $f(x)=\dfrac{x^2+\cos x}{x^2}$ pour $x\neq 0$.
	Montrer que $\displaystyle\lim_{x\to+\infty}f(x)=1$.
	
	$f(x)=1+\dfrac{\cos x}{x^2}$. Or $\dfrac{-1}{x^2}\leq\dfrac{\cos x}{x^2}\leq\dfrac{1}{x^2}\to 0$.
	
	Donc $\displaystyle\lim_{x\to+\infty}f(x)=1+0=1$.
\end{exemple}

% ----------------------------------------------------------------
\section{Limites et asymptotes}
% ----------------------------------------------------------------

\begin{definition}
	Soit $\mathcal{C}$ la courbe représentative de $f$.
	\begin{itemize}
		\item \textbf{Asymptote horizontale} $y=\ell$ :
		$\displaystyle\lim_{x\to\pm\infty}f(x)=\ell$.
		\item \textbf{Asymptote verticale} $x=a$ :
		$\displaystyle\lim_{x\to a}f(x)=\pm\infty$.
		\item \textbf{Asymptote oblique} $y=ax+b$ :
		$\displaystyle\lim_{x\to\pm\infty}[f(x)-(ax+b)]=0$.
	\end{itemize}
\end{definition}

\begin{methode}
	Pour chercher une asymptote oblique $y=ax+b$ :
	\begin{enumerate}
		\item Calculer $a=\displaystyle\lim_{x\to\pm\infty}\frac{f(x)}{x}$.
		\item Calculer $b=\displaystyle\lim_{x\to\pm\infty}[f(x)-ax]$.
		\item Si $a$ et $b$ sont finis, $y=ax+b$ est asymptote oblique.
	\end{enumerate}
\end{methode}

\begin{exemple}
	Soit $f(x)=\dfrac{x^2+2x-1}{x-1}$. Trouver les asymptotes.
	
	\textbf{Asymptote verticale :}
	$\displaystyle\lim_{x\to 1}(x-1)=0$ et $x^2+2x-1\to 2\neq 0$,
	donc $x=1$ est asymptote verticale.
	
	\textbf{Asymptote oblique :}
	$a=\displaystyle\lim_{x\to\pm\infty}\frac{x^2+2x-1}{x(x-1)}=1$.
	
	$f(x)-x=\dfrac{x^2+2x-1}{x-1}-x=\dfrac{x^2+2x-1-x(x-1)}{x-1}
	=\dfrac{3x-1}{x-1}\to 3$.
	
	Donc $y=x+3$ est asymptote oblique.
	
	\textbf{Position :} $f(x)-(x+3)=\dfrac{3x-1-3(x-1)}{x-1}=\dfrac{2}{x-1}$.
	Signe de $\dfrac{2}{x-1}$ : positif si $x>1$, négatif si $x<1$.
	La courbe est \textbf{au-dessus} de $\Delta$ pour $x>1$,
	\textbf{en dessous} pour $x<1$.
\end{exemple}

% ----------------------------------------------------------------
\section{Exercices résolus — Formes indéterminées}
% ----------------------------------------------------------------

\exercice{\etoilepleine\etoilevide\etoilevide}
Calculer les limites suivantes :
\[
\text{a)}\;\lim_{x\to+\infty}\frac{4x^3-x+2}{2x^3+5}
\qquad
\text{b)}\;\lim_{x\to 3}\frac{x^2-9}{x-3}
\qquad
\text{c)}\;\lim_{x\to-\infty}\frac{x^2-1}{x+1}
\]

\solution

\textbf{a)} FI $\dfrac{\infty}{\infty}$. Terme dominant $x^3$ :
\[
\lim_{x\to+\infty}\frac{4x^3(1-\frac{1}{4x^2}+\frac{1}{2x^3})}
{2x^3(1+\frac{5}{2x^3})} = \frac{4}{2} = 2.
\]

\textbf{b)} FI $\dfrac{0}{0}$. Factoriser :
$\dfrac{x^2-9}{x-3}=\dfrac{(x-3)(x+3)}{x-3}=x+3\to 6$.

\textbf{c)} FI $\dfrac{\infty}{\infty}$. Factoriser :
$\dfrac{x^2-1}{x+1}=\dfrac{(x-1)(x+1)}{x+1}=x-1\to-\infty$.

\bigskip
\exercice{\etoilepleine\etoilepleine\etoilevide}
Calculer :
\[
\text{a)}\;\lim_{x\to+\infty}\!\left(\sqrt{x^2+x}-x\right)
\qquad
\text{b)}\;\lim_{x\to 4}\frac{\sqrt{x}-2}{x-4}
\qquad
\text{c)}\;\lim_{x\to+\infty}\frac{x+\sin x}{2x+1}
\]

\solution

\textbf{a)} FI $+\infty-\infty$. Conjugué :
\[
\sqrt{x^2+x}-x
= \frac{x^2+x-x^2}{\sqrt{x^2+x}+x}
= \frac{x}{\sqrt{x^2+x}+x}
= \frac{x}{x(\sqrt{1+\frac{1}{x}}+1)}
= \frac{1}{\sqrt{1+\frac{1}{x}}+1}
\to \frac{1}{2}.
\]

\textbf{b)} FI $\dfrac{0}{0}$. Conjugué au numérateur :
\[
\frac{\sqrt{x}-2}{x-4}
= \frac{\sqrt{x}-2}{(\sqrt{x}-2)(\sqrt{x}+2)}
= \frac{1}{\sqrt{x}+2}
\to \frac{1}{4}.
\]

\textbf{c)} Gendarmes : $\dfrac{x-1}{2x+1}\leq\dfrac{x+\sin x}{2x+1}\leq\dfrac{x+1}{2x+1}$.
Les deux membres tendent vers $\dfrac{1}{2}$, donc la limite est $\dfrac{1}{2}$.

\bigskip
\exercice{\etoilepleine\etoilepleine\etoilepleine}
Soit $f(x)=\dfrac{2x^2-3x+1}{x-1}$.
\begin{enumerate}
	\item Simplifier $f(x)$ pour $x\neq 1$, puis calculer
	$\displaystyle\lim_{x\to 1}f(x)$.
	\item Trouver $a$ et $b$ tels que $f(x)=ax+b+\dfrac{c}{x-1}$
	(division euclidienne), puis donner l'asymptote oblique.
	\item Étudier la position de $\mathcal{C}_f$ par rapport à son asymptote oblique.
\end{enumerate}

\solution

\textbf{1.} $2x^2-3x+1=(2x-1)(x-1)$, donc pour $x\neq 1$ :
\[
f(x)=\frac{(2x-1)(x-1)}{x-1}=2x-1.
\]
Donc $\displaystyle\lim_{x\to 1}f(x)=1$.

\textbf{2.} On effectue la division de $2x^2-3x+1$ par $x-1$ :
\[
2x^2-3x+1 = (x-1)(2x-1)+0.
\]
Donc $f(x)=2x-1$ (pas de reste) : $a=2$, $b=-1$, $c=0$.

Asymptote oblique : $y=2x-1$ (la courbe \textbf{est} la droite pour $x\neq 1$).

\textbf{3.} $f(x)-(2x-1)=0$ pour tout $x\neq 1$ :
la courbe est confondue avec son asymptote oblique (sauf en $x=1$).

% ----------------------------------------------------------------
\section{Exercice type BAC}
% ----------------------------------------------------------------

\exercice{\etoilepleine\etoilepleine\etoilepleine}

Soit $f$ la fonction définie sur $\R\setminus\{1\}$ par :
\[
f(x) = \frac{x^3 - 3x + 2}{x^2 - 1}.
\]
On note $\mathcal{C}$ sa courbe représentative.

\medskip
\textbf{Partie A : Étude des limites et asymptotes}
\begin{enumerate}
	\item Calculer $\displaystyle\lim_{x\to+\infty}f(x)$,
	$\displaystyle\lim_{x\to-\infty}f(x)$.
	En déduire une asymptote oblique.
	\item Étudier les limites de $f$ en $x=1$ (à gauche et à droite).
	En déduire la nature de la droite $x=1$.
	\item Calculer $\displaystyle\lim_{x\to -1}f(x)$.
	La droite $x=-1$ est-elle asymptote verticale ?
\end{enumerate}

\textbf{Partie B : Position par rapport aux asymptotes}
\begin{enumerate}
	\setcounter{enumi}{3}
	\item Déterminer $a$, $b$, $c$ tels que
	$f(x)=ax+b+\dfrac{cx+d}{x^2-1}$ (division euclidienne).
	\item Étudier la position de $\mathcal{C}$ par rapport à l'asymptote oblique.
\end{enumerate}

\solution

\textbf{1.} FI $\dfrac{\infty}{\infty}$ :
$\displaystyle\lim_{x\to\pm\infty}\frac{x^3}{x^2}=\lim_{x\to\pm\infty}x$.
Donc $\displaystyle\lim_{x\to+\infty}f(x)=+\infty$ et
$\displaystyle\lim_{x\to-\infty}f(x)=-\infty$.

Cherchons $a$ et $b$ :
$a=\displaystyle\lim_{x\to+\infty}\frac{f(x)}{x}
=\lim_{x\to+\infty}\frac{x^3-3x+2}{x(x^2-1)}
=\lim_{x\to+\infty}\frac{x^3}{x^3}=1$.

$b=\displaystyle\lim_{x\to+\infty}[f(x)-x]
=\lim_{x\to+\infty}\frac{x^3-3x+2-x(x^2-1)}{x^2-1}
=\lim_{x\to+\infty}\frac{-2x+2}{x^2-1}
=\lim_{x\to+\infty}\frac{-2}{x}=0$.

Donc $\Delta : y=x$ est asymptote oblique en $\pm\infty$.

\textbf{2.} En $x=1$ : $x^2-1\to 0$ et
$x^3-3x+2\big|_{x=1}=1-3+2=0$.
FI $\dfrac{0}{0}$. On factorise :
$x^3-3x+2=(x-1)^2(x+2)$ et $x^2-1=(x-1)(x+1)$.
\[
f(x)=\frac{(x-1)^2(x+2)}{(x-1)(x+1)}=\frac{(x-1)(x+2)}{x+1}
\quad (x\neq\pm 1).
\]
Donc $\displaystyle\lim_{x\to 1}f(x)=\frac{0\cdot3}{2}=0$.
La droite $x=1$ n'est \textbf{pas} une asymptote verticale
($f$ admet une limite finie en 1 : c'est un trou).

\textbf{3.} En $x=-1$ :
$\displaystyle\lim_{x\to-1}\frac{(x-1)(x+2)}{x+1}$.
Numérateur : $(-2)(1)=-2\neq 0$. Dénominateur $\to 0$.

Par valeurs à droite ($x\to-1^+$) : $\dfrac{-2}{0^+}=-\infty$.
Par valeurs à gauche ($x\to-1^-$) : $\dfrac{-2}{0^-}=+\infty$.

Donc $x=-1$ est bien une \textbf{asymptote verticale}.

\textbf{4.} Division euclidienne de $x^3-3x+2$ par $x^2-1$ :
\[
x^3-3x+2 = x(x^2-1)+(-2x+2).
\]
Donc $f(x)=x+\dfrac{-2x+2}{x^2-1}=x+\dfrac{-2(x-1)}{(x-1)(x+1)}
=x-\dfrac{2}{x+1}$ (pour $x\neq 1$).

\textbf{5.} $f(x)-x=-\dfrac{2}{x+1}$.

Pour $x>-1$ (et $x\neq 1$) : $f(x)-x<0$ : $\mathcal{C}$ est
\textbf{en dessous} de $\Delta$.

Pour $x<-1$ : $f(x)-x>0$ : $\mathcal{C}$ est
\textbf{au-dessus} de $\Delta$.
